\section{Experiments and Results}\label{sec:experiments}
The controlled benchmark tests the class-probability mechanism, followed by independent mapping comparisons and semantic stress tests. The public check provides auxiliary evidence from measured metadata. Scorer coefficients and gate rules remain fixed throughout these comparisons.

\subsection{Controlled descriptor benchmark and evaluation protocol}
\subsubsection{Synthetic multi-receiver scenes}
A nominal scene places four receivers on a randomly rotated circle of radius 2000 m. It contains $E=14$ hidden \emph{transmission instances}, indexed $e=0,\ldots,13$, each with its own time, position, carrier, bandwidth, duration, and anonymous family label. The 14 bursts are not 14 classes, named protocols, reports, or persistent emitters. Independent detection gives zero to four reports per burst; at probability 0.8, the expected genuine total is $14\times4\times0.8=44.8$, before clutter.

Initial positions are drawn in a $3000\times3000$ m square and start times in a $120\,\mu$s window. Carrier and bandwidth are drawn independently before near-collision overrides:
\begin{equation}
 \begin{aligned}
 f_e^{(0)}&=100\,\mathrm{MHz}+m_e(1.8\,\mathrm{kHz})+\eta_e,
 &&m_e\sim\operatorname{Unif}\{0,\ldots,5\},\\
 \eta_e&\sim\mathcal N(0,(180\,\mathrm{Hz})^2),
 &&B_e^{(0)}\sim\operatorname{Unif}[20,60]\,\mathrm{kHz}.
 \end{aligned}
 \label{eq:carrier_generator}
\end{equation}
The 1.8 kHz carrier spacing and 20--60 kHz bandwidth deliberately create overlapping descriptors: six centers span only 9 kHz. The 180 Hz event perturbation $\eta_e$ is shared across reports and cancels from true-pair differences. It is separate from receiver residuals of 250 Hz and report errors of 120--360 Hz (Table~\ref{tab:configuration}); the 180 Hz value is not a waveform spectral width.

Initial durations are uniform over 20--80 $\mu$s. Each event after the first has probability 0.25 of copying an earlier event's time, carrier, bandwidth, and duration with small perturbations, while retaining independent position and class. These overrides can move parameters beyond the initial ranges.

The collisions create descriptor ambiguity without simulating interference, packet detection, or overlapping-waveform demodulation. All 14 first-scene events and detecting receivers are recorded in \texttt{results/event\_ledger\_r00\_test000.csv}.

\begin{table}[!t]
\begin{adjustwidth}{-\extralength}{0cm}
\caption{Nominal generator and inference settings. Receiver index $s$ starts at zero. Standard deviations describe measurement errors unless marked as shared residual biases. Supplementary Sections S1--S2 give the generation order and stress changes.}\label{tab:configuration}
\small\renewcommand{\arraystretch}{1.12}
\begin{tabularx}{\linewidth}{@{}p{5.1cm}Y@{}}\toprule
Quantity & Configuration \\\midrule
Receivers; transmission instances; classes & $S=4$; $E=14$; $K=4$ anonymous signal-family labels \\
Nominal carrier-center spacing & 1.8 kHz; six centers around 100 MHz, not channel bandwidth \\
Event carrier perturbation s.d. & 180 Hz, shared across reports of one event before collision overrides \\
Initial event bandwidth; duration & Uniform 20--60 kHz; uniform 20--80 $\mu$s, before collision overrides \\
Detection; clutter & Detection probability 0.8; independent Poisson mean 0.8 per receiver \\
Timing; carrier measurement s.d. & $(0.6+0.4s)\,\mu$s; $(0.12+0.08s)$ kHz \\
Bandwidth; duration measurement s.d. & $(1.5+0.3s)$ kHz; $(1+0.3s)\,\mu$s \\
Scene-shared receiver residuals & Clock s.d. $0.8\,\mu$s; oscillator s.d. 0.25 kHz; correctly declared nominally \\
Semantic corruption; missingness & Wrong-class probability 0.12; missing-vector probability 0.10 \\
Reported class confidence & Correct: uniform $[0.60,0.95]$; incorrect: uniform $[0.45,0.90]$ \\
Linear scorer; calibration & $C_{\rm reg}=1$, maximum 500 iterations; $T\in[0.1,10]$ \\
Boosted comparator & 100 iterations; 15 maximum leaf nodes; learning rate 0.08; $L_2=1$ \\
Gate; greedy numerical cutoff & $k=3.5$; positive merge gain exceeding $10^{-10}$ \\\bottomrule
\end{tabularx}\end{adjustwidth}
\end{table}

Detection probability is 0.8 with independent clutter. Clock and oscillator residuals are scene-shared, whereas detector errors are drawn per report. Native frequency units and class permutations are applied before shuffling.

Generated class vectors give the designated family confidence in $[0.60,0.95]$ when correct and $[0.45,0.90]$ when wrong, distributing remaining mass equally. Missing vectors are uniform (Table~\ref{tab:configuration}). These are synthetic inputs, not trained waveform-classifier outputs; Section~\ref{sec:public_data} treats measured records separately.

\subsubsection{Scene splits and reproducible implementation}
Ten replications independently regenerate the scenes and refit scorers. Replication $r\in\{0,\ldots,9\}$ uses seed $10^6+10^5r$, split offsets 0, 10,000, and 20,000, and the scene index. Splitting 120/40/80 scenes before pairing prevents event or receiver-residual leakage across splits \cite{kapoor}.

\begin{table}[!t]\caption{Nominal data totals over ten replications. Positive pairs represent a common true event; all listed pairs are cross-receiver.}\label{tab:splits}
\centering\small\renewcommand{\arraystretch}{1.12}
\begin{tabular}{lrrrr}\toprule
Split & Scenes & Observations & Pairs & Positive pairs \\\midrule
Training & 1200 & 57,227 & 1,023,517 & 63,897 \\
Calibration & 400 & 19,113 & 342,476 & 21,395 \\
Test & 800 & 38,451 & 693,023 & 43,151 \\\bottomrule
\end{tabular}\end{table}

Table~\ref{tab:splits} gives observed counts. The implementation uses NumPy, SciPy, and scikit-learn \cite{sklearn,scipy}. Original runs used Python 3.13.5, NumPy 2.3.5, SciPy 1.17.0, scikit-learn 1.8.0, and pandas 2.2.3 on Linux with an AMD EPYC 9V74 processor. Added experiments record a Windows environment. Numerical libraries use one thread without a GPU; original timing curves exclude generation, features, fitting, and scoring.

\subsubsection{Comparators, ablations, and metrics}
We compare timing-only, frequency-only, semantics-only, physical, and six-input linear or boosted scorers. Raw-frequency ablation removes conversion only from the frequency \emph{feature}, retaining the canonical gate and other inputs. Separate variants remove entropy weighting, gating, or temperature scaling.

All decoders share observations and pair evidence. Components join admitted edges above 0.5. Pairwise assignment allows zero-gain dummies and then unions receiver-pair matchings, without optimizing a global multipartite assignment.

Constrained scorers use Algorithm~\ref{alg:greedy}. Ablations remove exclusivity or complete cross-pair admission separately. Without exclusivity, unscored same-receiver edges have zero weight. Stress tests retain nominal fitted parameters.

AUROC and non-interpolated average precision (AP) assess discrimination; Brier error and clipped negative log-likelihood (NLL) assess probabilities \cite{silvafilho2023,miller2024}. Pair tables include rejected pairs, whereas reliability diagrams use admitted pairs in ten equal-width bins.

Partition metrics use \emph{all unordered observation pairs}, so same-receiver false links are penalized. Event-pair F1 is
\begin{equation}
 F_1=\frac{2\,\mathrm{TP}}{2\,\mathrm{TP}+\mathrm{FP}+\mathrm{FN}},\label{eq:f1}
\end{equation}
TP, FP, and FN count common, predicted-only, and true-only co-memberships. We also report ARI, $B^3$ F1, exact visible-event recovery, and clutter attachment \cite{hassan2024,miller2024}. Exact recovery requires an event seen at least twice to be recovered without additions or omissions. Receiver violations use non-singleton clusters, while gate violations use predicted within-cluster cross-receiver edges (Supplementary Section S3).

Pair metrics pool 80 scenes per replication; partition metrics first average scene scores. Ten-replication means use 95\% intervals $\bar x\pm t_{0.975,9}s_x/\sqrt{10}$, with Student $t$ on nine degrees of freedom and replication standard deviation $s_x$. The F1 estimand is $\frac1{10}\sum_{r=1}^{10}\frac1{80}\sum_{b=1}^{80}F_{1,rb}$. Joint constrained counts are TP $=30{,}063$, FP $=3589$, and FN $=13{,}088$; TP $+$ FN equals the 43,151 test positives (Table~\ref{tab:splits}). Micro F1 is 0.7829 versus scene-macro F1 of 0.7798. Supplementary Section S14 gives all counts; method differences are paired before interval calculation.

The 693,023 dependent pairs are not independent replications. Intervals describe generation and refitting under this model, not variation across field deployments.

\subsection{Representation and known-vocabulary alignment checks}\label{sec:alignment_audit}
The mapping audit reuses 800 scenes and ten frozen joint-linear scorers with fixed gates and decoding. Correctly mapped, bypassed native, and uniform class vectors change the input only, without new scenes or fitting.

\begin{table}[!t]\begin{adjustwidth}{-\extralength}{0cm}
\caption{Event alignment under supplied class maps. All rows use the same frozen joint scorer, physical gate, decoder, and scenes. F1 scores event matches. It does not measure recovery of an unknown vocabulary. Intervals use ten replication-level means.}\label{tab:semantic_alignment}
\small\renewcommand{\arraystretch}{1.12}\begin{tabular*}{\linewidth}{@{\extracolsep{\fill}}lrrr@{}}\toprule
Input treatment & Event-pair F1 [95\% CI] & Receiver conflict & Gate violation \\\midrule
Declared class mapping applied & 0.7798 [0.7713, 0.7883] & 0.00\% & 0.00\% \\
Native columns treated as common & 0.4521 [0.4366, 0.4676] & 0.00\% & 0.00\% \\
Class evidence set to uniform & 0.6438 [0.6290, 0.6587] & 0.00\% & 0.00\% \\
\bottomrule\end{tabular*}

\end{adjustwidth}\end{table}

Mapped F1 is 0.7798 versus 0.4521 for bypass, a paired difference of 0.3277 with 95\% interval $[0.3162,0.3392]$. Uniform vectors give 0.6438, compared with 0.6813 from separately refitted physical scoring. Input removal and refitting answer different questions. All conditions remain feasible, so structural checks do not correct class meaning.

Matched unit and class-order changes preserve all 800 partitions with zero feature or probability difference, checking Equation~\eqref{eq:class_mapping}. The 800 re-encodings reuse scenes and do not test independent accuracy or unknown, partial, or hierarchical maps. Outputs are retained in \texttt{results/semantic\_alignment/}.

\subsection{Pair discrimination and confidence quality}
Class evidence raises calibrated linear AP from 0.6890 to 0.7979 (Table~\ref{tab:pair}), a paired gain of 0.1089 with 95\% interval $[0.1015,0.1164]$. It resolves some physical ambiguity that moderate single-modality AUROC obscures. Boosting reaches AP 0.8343, indicating useful nonlinear interactions.

\begin{table}[!t]\begin{adjustwidth}{-\extralength}{0cm}
\caption{Pair results over ten independently fitted replications. Higher AUROC and AP are better. Lower Brier error and NLL are better. $T$ means temperature scaling. AP intervals use variation across replications.}\label{tab:pair}
\small\renewcommand{\arraystretch}{1.12}\begin{tabularx}{\linewidth}{@{}Yrrrr@{}}
\toprule
Method & AUROC & AP [95\% CI] & Brier & NLL \\
\midrule
Timing only & 0.9135 & 0.2864 [0.2799, 0.2928] & 0.0469 & 0.1435 \\
Frequency only & 0.8995 & 0.2550 [0.2476, 0.2624] & 0.0502 & 0.1611 \\
Semantics only & 0.7785 & 0.1607 [0.1583, 0.1630] & 0.0563 & 0.2119 \\
Physical, ungated & 0.9831 & 0.6885 [0.6744, 0.7027] & 0.0253 & 0.0741 \\
Joint linear, ungated & 0.9885 & 0.7963 [0.7878, 0.8048] & 0.0211 & 0.0644 \\
Boosted joint, ungated & 0.9910 & 0.8346 [0.8278, 0.8414] & 0.0183 & 0.0562 \\
Raw-frequency feature & 0.9868 & 0.7757 [0.7665, 0.7850] & 0.0227 & 0.0696 \\
Unweighted semantics, ungated & 0.9899 & 0.8202 [0.8120, 0.8284] & 0.0198 & 0.0611 \\
Joint linear, gated & 0.9884 & 0.7979 [0.7893, 0.8065] & 0.0209 & 0.0638 \\
Joint linear, gated + $T$ & 0.9884 & 0.7979 [0.7893, 0.8065] & 0.0208 & 0.0637 \\
Physical, gated + $T$ & 0.9829 & 0.6890 [0.6747, 0.7032] & 0.0250 & 0.0734 \\
Boosted joint, gated + $T$ & 0.9908 & 0.8343 [0.8274, 0.8412] & 0.0183 & 0.0569 \\
\bottomrule
\end{tabularx}

\end{adjustwidth}\end{table}

\pairedfig{results_pr}{All cross-receiver test pairs.}{results_reliability}{Admitted pairs of the joint linear model.}{Diagnostics for held-out replication 0. (a) Precision--recall curves compare physical, joint linear, and boosted scores with gating and calibration. (b) Ten-bin confidence summaries and their post-minus-pre differences. Horizontal displacement $\pm0.006$ improves legibility without altering values. Tables provide uncertainty across replications.}{fig:pair_results}

Figure~\ref{fig:pair_results}a compares discrimination; Figure~\ref{fig:pair_results}b evaluates admitted-pair confidence, matching the temperature-fitting population. Scaling changes paired Brier error by $-2.88\times10^{-5}$ with interval $[-5.60\times10^{-5},-1.55\times10^{-6}]$: statistically nonzero but numerically modest.

Positive rescaling preserves unsaturated ranking. All 800 executable nominal partitions remain unchanged, consistent with Proposition~\ref{prop:temperature}; calibration changes confidence without clustering gain.

Removing concentration weighting raises ungated AP from 0.7963 to 0.8202, showing that diffuse compatibility can remain useful. Raw-frequency AP falls to 0.7757, with other features and the canonical gate partly compensating. Entropy weights are not a universal reliability estimate, and this frequency ablation does not remove the full observation contract.

\subsection{Event partitions and the role of explicit constraints}
Joint constrained F1 is 0.7798 versus physical-only F1 of 0.6813 (Table~\ref{tab:cluster}; Figure~\ref{fig:partition_results}), a paired gain of 0.0985 with interval $[0.0887,0.1083]$. Components yield F1 0.6764, recall 0.8985, and precision 0.5703 because unchecked chains can create oversized or receiver-conflicting groups.

\begin{table}[!t]\begin{adjustwidth}{-\extralength}{0cm}
\caption{Scorer--decoder combinations on 800 test scenes. Decoders share each scorer's probabilities and gates; F1 intervals use ten refitted replications. Receiver conflicts are counted per non-singleton cluster and gate violations per predicted within-cluster cross-receiver edge.}\label{tab:cluster}
\small\renewcommand{\arraystretch}{1.10}\begin{tabularx}{\linewidth}{@{}llYrrrrr@{}}\toprule
Scorer & Decoder & Event-pair F1 [95\% CI] & ARI & $B^3$ F1 & Exact & Rx conflict & Gate violation \\\midrule
Physical linear & Components & 0.6594 [0.6472, 0.6717] & 0.6369 & 0.8432 & 0.4428 & 22.2\% & 0.5\% \\
Physical linear & Assignment + union & 0.7145 [0.7038, 0.7252] & 0.6981 & 0.8543 & 0.4981 & 12.7\% & 0.2\% \\
Physical linear & Constrained & 0.6813 [0.6658, 0.6968] & 0.6683 & 0.8129 & 0.4072 & 0.0\% & 0.0\% \\
Joint linear & Components & 0.6764 [0.6641, 0.6887] & 0.6549 & 0.8540 & 0.4909 & 21.1\% & 0.7\% \\
Joint linear & Assignment + union & 0.7797 [0.7741, 0.7853] & 0.7671 & 0.8898 & 0.5973 & 10.5\% & 0.1\% \\
Joint linear & Constrained & 0.7798 [0.7713, 0.7883] & 0.7704 & 0.8677 & 0.5375 & 0.0\% & 0.0\% \\
Unweighted joint & Components & 0.6889 [0.6780, 0.6998] & 0.6684 & 0.8597 & 0.5067 & 19.7\% & 0.6\% \\
Unweighted joint & Assignment + union & 0.8025 [0.7974, 0.8075] & 0.7915 & 0.8993 & 0.6238 & 8.6\% & 0.1\% \\
Unweighted joint & Constrained & 0.7983 [0.7892, 0.8074] & 0.7897 & 0.8784 & 0.5652 & 0.0\% & 0.0\% \\
Boosted joint & Components & 0.7016 [0.6888, 0.7145] & 0.6819 & 0.8676 & 0.5354 & 18.6\% & 0.9\% \\
Boosted joint & Assignment + union & 0.8130 [0.8063, 0.8196] & 0.8024 & 0.9072 & 0.6506 & 8.6\% & 0.2\% \\
Boosted joint & Constrained & 0.8238 [0.8181, 0.8295] & 0.8160 & 0.8928 & 0.6189 & 0.0\% & 0.0\% \\
\bottomrule\end{tabularx}

\end{adjustwidth}\end{table}

\widefig{results_partition}{Event-pair F1 for each scorer and decoder. Points show means of ten replication-level averages. Bars show 95\% $t$ intervals. All methods use the same held-out scenes. Unweighted semantics and boosted scoring both give higher F1 than the joint constrained reference.}{fig:partition_results}

Receiver conflicts affect 21.14\% of non-singleton component clusters and 10.49\% of assignment clusters. Algorithm~\ref{alg:greedy} prevents receiver and gate violations, but its precision 0.8932 and recall 0.6972 leave many true pairs unjoined.

Assignment gives similar F1 (0.7797) and higher exact recovery (0.5973 versus 0.5375), while producing some infeasible groups. Accuracy and structural validity must therefore be assessed together.

Table~\ref{tab:cluster} crosses four scorers and three decoders. Boosted F1 is 0.7016 for components, 0.8130 for assignment, and 0.8238 for constraints. The constrained-minus-assignment gain is 0.0108 with 95\% interval $[0.0020,0.0196]$, while assignment retains 8.57\% receiver conflicts.

Physical-only constrained F1 is 0.6813 versus assignment F1 of 0.7145, a difference of $-0.0332$ with interval $[-0.0439,-0.0225]$. Joint and unweighted intervals include zero. Constraint accuracy depends on evidence, while feasibility is common to all constrained scorers. The 9600 outcomes and frozen probabilities reproduce the retained results.

Removing exclusivity lowers F1 to 0.7361 and repeats receivers. Removing complete admission gives F1 0.7796 but permits rejected within-group edges, making this check principally a nominal feasibility safeguard.

Unweighted and boosted scoring give F1 0.7983 and 0.8238. The former motivates the simpler linear default; weighted scoring remains the historical control. Independent validation retains all three semantic scorers for event and missing-map tests.

\subsection{Admission sensitivity and distribution shifts}
\subsubsection{Gate width and irreversible pair removal}
The frozen-scorer sweep uses the first 20 scenes per replication (Table~\ref{tab:gate}). At $k=1$, true-pair recall is 68.37\% and F1 is 0.5975; later semantics cannot restore rejected edges. At $k=3.5$, recall is 0.9997 with 17.72\% of candidates retained. Increasing $k$ to 5 adds competition with little gain. At $k=3.5$, F1 0.7822 differs from 0.7798 because the sweep uses 20 rather than 80 scenes.

\begin{table}[!t]\begin{adjustwidth}{-\extralength}{0cm}
\caption{Gate-width test on 200 nominal scenes: ten replications times the first 20 test scenes. Retained pair fraction counts admitted cross-receiver pairs. True-pair recall uses only pairs from a common true event.}\label{tab:gate}
\small\renewcommand{\arraystretch}{1.12}\begin{tabularx}{\linewidth}{@{}Yrrr@{}}
\toprule
$k$ & True-pair recall & Retained pair fraction & Event-pair F1 [95\% CI] \\
\midrule
1 & 0.6837 & 0.0885 & 0.5975 [0.5806, 0.6145] \\
2 & 0.9585 & 0.1354 & 0.7569 [0.7394, 0.7744] \\
3 & 0.9979 & 0.1621 & 0.7808 [0.7620, 0.7996] \\
3.5 & 0.9997 & 0.1772 & 0.7822 [0.7635, 0.8009] \\
5 & 1.0000 & 0.2254 & 0.7825 [0.7640, 0.8010] \\
\bottomrule
\end{tabularx}

\end{adjustwidth}\end{table}

\subsubsection{Incorrect semantics and misspecified receiver uncertainty}
The frozen-model grid has 34 settings: 12 named conditions, six semantic-error probabilities, 12 actual-versus-declared clock combinations, and four scene sizes. Each uses 20 scenes per replication with seeds disjoint from nominal splits. Reusing stress schedules couples comparable draws; the 34 settings are controlled interventions, not separate measured datasets.

\begin{table}[!t]\begin{adjustwidth}{-\extralength}{0cm}
\caption{Named stress tests with 200 scenes per condition. The separated first row is the stress-seed reference. It differs from the nominal test subset in Table~\ref{tab:gate}. The semantic-error row uses error probability 0.45. All other settings follow the stated intervention. The scorer, scaler, gate multiplier, and temperature stay fixed. F1 is event-pair F1. The last column is true-pair gate recall.}\label{tab:stress}
\small\renewcommand{\arraystretch}{1.12}\begin{tabularx}{\linewidth}{@{}Yrrr@{}}
\toprule
Test condition & Physical event F1 & Joint event F1 [95\% CI] & Gate recall \\
\midrule
\multicolumn{4}{l}{\textit{Reference condition on the stress seeds}}\\
Nominal stress reference & 0.6692 & 0.7758 [0.7491, 0.8024] & 0.9996 \\
\midrule
\multicolumn{4}{l}{\textit{Interventions on the same seed schedule}}\\
28 events & 0.6754 & 0.7601 [0.7478, 0.7724] & 0.9997 \\
Detection probability 0.55 & 0.6770 & 0.7652 [0.7472, 0.7831] & 0.9993 \\
Clutter mean 2.5 / receiver & 0.6881 & 0.7862 [0.7754, 0.7970] & 0.9997 \\
Semantic error probability 0.45 & 0.6727 & 0.6308 [0.6093, 0.6523] & 0.9992 \\
Missing semantics probability 0.65 & 0.6692 & 0.6612 [0.6369, 0.6855] & 0.9996 \\
Clock s.d. $6\,\mu$s; declared $0.8\,\mu$s & 0.5303 & 0.6255 [0.5996, 0.6513] & 0.9235 \\
Frequency bias s.d. 1.5 kHz; declared 0.25 kHz & 0.3265 & 0.4110 [0.3845, 0.4375] & 0.5606 \\
Additional positive timing delays & 0.6377 & 0.7145 [0.6960, 0.7330] & 0.9549 \\
Variance-matched $t_3$ timing noise & 0.6885 & 0.7786 [0.7632, 0.7941] & 0.9993 \\
Fifth receiver, not used in training & 0.6743 & 0.7780 [0.7626, 0.7934] & 0.9998 \\
One confidently wrong receiver & 0.6692 & 0.5117 [0.4970, 0.5264] & 0.9996 \\
\bottomrule
\end{tabularx}

\end{adjustwidth}\end{table}

\pairedfig{results_semantic}{Increasing semantic-label error.}{results_clock_f1}{Increasing receiver clock residual.}{Frozen-model stress tests. Upper panels show ten-replication means and 95\% intervals; lower panels enlarge method differences on separate scales. Symbols shift by $\pm0.007$ in (a) and $\pm0.12\,\mu$s in (b), with trend lines at true values. Incorrect class confidence can harm joint scoring; correct clock declarations reduce but do not remove timing ambiguity.}{fig:stress_results}

The stress reference gives F1 0.7758 (Table~\ref{tab:stress}, first row). At semantic-error probability 0.45, joint F1 is 0.6308 versus physical F1 of 0.6727 (Table~\ref{tab:stress}). One confidently wrong receiver reduces joint F1 to 0.5117: its family is shifted to $(q_e+1)\bmod K$ with confidence 0.99, remaining mass 0.01, and missingness disabled. Equation~\eqref{eq:entropy} cannot detect these concentrated errors, defining a failure regime in Figure~\ref{fig:stress_results}a.

Actual oscillator s.d. 1.5 kHz versus declared 0.25 kHz reduces gate recall to 0.5606 and F1 to 0.4110. With clock s.d. $20\,\mu$s, declaring $0.8\,\mu$s gives F1 0.3136; matching the declaration raises it to 0.5486.

Gate recall rises from 0.4581 to 0.9996 (Figure~\ref{fig:runtime_results}a), but broader uncertainty also admits more competing matches. Restoring coverage alone does not restore accurate grouping.

Positive delays reduce gate coverage, while variance-matched heavy tails have a smaller effect here. The fifth-receiver test changes network size within the generator, not arbitrary receiver domains. Clutter and detection changes alter visibility and metric denominators, so higher scores do not establish a benefit from clutter or missed detections.

\FloatBarrier
\subsection{Computational cost, exact audit, and reproducibility}
\pairedfig{results_clock_gate}{Gate coverage under clock misspecification.}{results_runtime}{Measured dense-decoder runtime.}{Coverage and runtime diagnostics. (a) Fixed and matched uncertainty, with their recall difference below; symbols shift by $\pm0.12\,\mu$s only for legibility. (b) Physical and joint runtimes use equally scaled log axes and a percentage-difference strip. Event counts 8, 16, 32, and 64 yield different report counts. Timings exclude generation and scoring; ten-replication intervals are specific to this hardware.}{fig:runtime_results}

Figure~\ref{fig:runtime_results}b shows increasing dense-decoder work. Four sizes do not establish an asymptotic exponent; Proposition~\ref{prop:greedy} provides the $O(N^3)$ bound. The implementation retains dense matrices, so no sparse-implementation speedup is demonstrated.

Across 60 scenes with 11--19 reports, greedy and exact objectives agree in 47 cases. Relative gap $(J^*-J_g)/\max(|J^*|,10^{-9})$ averages 1.71\% and peaks at 16.82\%. Mean F1 is 0.7053 for greedy and 0.7182 for exact inference.

This comparison holds the \emph{same learned objective} fixed. It measures small-instance optimization loss, not larger-instance near-optimality or a monotone relation to true F1. Figure~\ref{fig:decoder} illustrates the greedy limit; detailed gaps and selective-pair risk--coverage curves remain in the supplement.

Audits regenerate all 2400 nominal scenes, verify splits and counts, recompute summaries and invariants, and compare predictions with fitted parameters. Stored outputs retain 6400 original rows and the separate 9600-row factorial, whose accuracy and feasibility are recomputed from saved labels.

Scientific outputs repeat within floating-point tolerance, while runtimes depend on execution conditions. Subsequent assumption tests remain descriptor-based.
\subsection{Single-receiver exploratory mapping probe}\label{sec:partial_probe}
Table~\ref{tab:partial_probe} reports an exploratory test that hides three entries at one receiver. One existing calibration scene recovers 26/30 entries across ten repetitions; five or 40 scenes recover 30/30. Because these descriptors were previously used for temperature fitting, they do not form an independent adaptation set. Supplementary Section S9 gives the protocol.

The measured-input probe hides an artificial swap of decoded SF7/SF12 meanings at one LoED gateway. Adaptation uses 3 August. Testing uses 2 September. Recovering the two entries restores reference-packet F1 from 0.1270 to 0.9961. This is code registration in decoded metadata. The experiment contains no measured heterogeneous class posteriors. Sequential and alternating estimates tie. The independent multi-receiver experiment below provides the stronger controlled mapping test.
\begin{table}[!t]\begin{adjustwidth}{-\extralength}{0cm}
\caption{Optional missing-map test with the original scores and decoder. Controlled results average ten repetitions. Public results use one later test capture. Mapping accuracy counts only hidden entries. Controlled rows use event-pair F1. Public rows use reference-packet F1.}\label{tab:partial_probe}
\small\setlength{\tabcolsep}{5pt}\renewcommand{\arraystretch}{1.10}
\begin{tabularx}{\linewidth}{@{}lYrr@{}}\toprule
Probe / adaptation & Method & Map (\%) & Pair F1\\\midrule
Controlled / 1 scene & Unresolved completion & 33.3 & 0.7026\\
Controlled / 1 scene & Sequential map fit & 86.7 & 0.7639\\
Controlled / 1 scene & Map--event feedback & 86.7 & 0.7639\\
Controlled / 40 scenes & Sequential map fit & 100.0 & 0.7798\\
Controlled / 40 scenes & Map--event feedback & 100.0 & 0.7798\\
Controlled / 40 scenes & Complete-map reference & 100.0 & 0.7798\\
\midrule
Public / 1 capture & Unresolved completion & 0.0 & 0.1270\\
Public / 1 capture & Sequential map fit & 100.0 & 0.9961\\
Public / 1 capture & Map--event feedback & 100.0 & 0.9961\\
Public / 1 capture & Complete-map reference & 100.0 & 0.9961\\
\bottomrule\end{tabularx}
\par\vspace{3pt}\footnotesize The controlled test hides three entries at one receiver. The public test hides only SF7/SF12 at one gateway. Five-scene controlled adaptation also reaches 100\% and F1 0.7798. The public test has 8423 reports. The original public study pools 23,677 reports. Feedback does not improve on sequential fitting.
\end{adjustwidth}\end{table}


\subsection{Map robustness and multi-receiver completion}\label{sec:map_revision}
\subsubsection{Incorrect relative mappings and shared renaming}
The 800-scene wrong-map audit freezes the joint scorers and changes 0, 2, 4, 8, 12, or 16 of 16 entries through bijections. Three interventions per budget are averaged within replication. At 25\% corruption, F1 falls from 0.7798 to 0.6422; at 50\%, it is 0.5474 (Figure~\ref{fig:review_mapping}a). All groups remain feasible, showing that constraints cannot repair class meaning.

The curve is nonmonotone because random permutations can retain relative agreement. A shared renaming changes every absolute entry but preserves all 800 partitions. Event F1 therefore cannot establish absolute naming accuracy.

\subsubsection{Independent adaptation and conditional map optimization}
Each replication supplies 40 adaptation scenes disjoint from all original splits and stress seeds, without refitting. One to three maps are hidden while receiver 0 stays known; unknown receivers receive zero to two anchors. Fixed prefix budgets 1, 5, and 40 are all reported, without test-selected stopping. The 27 configurations use ten seeds and reuse the initial 800 tests. This analysis followed inspection of those tests; the subsequent comparison uses independent tests.

For three unknown maps without anchors, one scene recovers 75/120 entries with F1 0.6787. Five scenes recover 120/120 and restore F1 0.7798 (Table~\ref{tab:review_mapping}), as do forty. Budgets two to four were untested, so five is only the first successful tested value. The ten map interventions do not constitute 120 independent Bernoulli trials or establish larger-vocabulary sample complexity.

Five-scene coordinate fitting recovers 60\% of entries with F1 0.7171, whereas exact fitting recovers all. In replication 0 with fixed groups, every single-map gain is nonpositive, but joint updating raises the objective from 201.0133 to 235.4834. Coupling between maps causes this trap despite complete source-domain enumeration. Exact sequential and alternating fitting tie in every tested configuration, with no added event-update gain.

\begin{table}[!t]\begin{adjustwidth}{-\extralength}{0cm}
\caption{Independent adaptation without target anchors. The reference map remains known. Counts sum ten interventions; F1 uses the same 800 tests. Coordinate updates one map at a time; sequential and alternating methods solve the exact conditional block. All receiver and gate violations are zero.}\label{tab:review_mapping}
\small\setlength{\tabcolsep}{3pt}\begin{tabularx}{\linewidth}{@{}rrrYrYr@{}}\toprule
Adaptation scenes & Unknown sources & Correct hidden entries & Sequential F1 & Feedback F1 [95\% CI] & Coordinate F1 & Rx / gate \\\midrule
1 & 1 & 35/40 & 0.7604 & 0.7604 [0.7280, 0.7929] & 0.7604 & 0 / 0 \\
1 & 2 & 61/80 & 0.7163 & 0.7163 [0.6628, 0.7697] & 0.6941 & 0 / 0 \\
1 & 3 & 75/120 & 0.6787 & 0.6787 [0.6272, 0.7301] & 0.6279 & 0 / 0 \\
5 & 1 & 40/40 & 0.7798 & 0.7798 [0.7713, 0.7883] & 0.7798 & 0 / 0 \\
5 & 2 & 80/80 & 0.7798 & 0.7798 [0.7713, 0.7883] & 0.7798 & 0 / 0 \\
5 & 3 & 120/120 & 0.7798 & 0.7798 [0.7713, 0.7883] & 0.7171 & 0 / 0 \\
40 & 1 & 40/40 & 0.7798 & 0.7798 [0.7713, 0.7883] & 0.7798 & 0 / 0 \\
40 & 2 & 80/80 & 0.7798 & 0.7798 [0.7713, 0.7883] & 0.7798 & 0 / 0 \\
40 & 3 & 120/120 & 0.7798 & 0.7798 [0.7713, 0.7883] & 0.7133 & 0 / 0 \\
\bottomrule\end{tabularx}

\end{adjustwidth}\end{table}

\widefig{review_mapping}{Mapping interventions. (a) Relative-map errors harm association, whereas shared renaming preserves it. (b,c) Missing fractions use all 16 entries, with the reference known and no target anchors. Five- and 40-scene curves coincide. F1 bars use replication means; hidden-entry fractions are descriptive. All violation rates are zero.}{fig:review_mapping}

Schema tests reject many-to-one maps, three-column coarse vectors, and a fifth unknown class with positive mass. This defines the implementation boundary, not open-set or hierarchical recovery. Supplementary Section S11 records anchors and maps.

\subsection{Independent comparison of missing-map estimators}\label{sec:deep_mapping}
Independent validation uses 400 adaptation and 800 test scenes, with ten replications of 40 and 80 scenes. Prefix budgets are 1, 5, and 40. Receiver 0 remains known; three unanchored four-class maps are hidden. All solvers share physical-only groups and frozen parameters. Reconstructing the boosted model from original training with its saved temperature reproduces all 800 partitions, with maximum probability difference below $4.5\times10^{-11}$.

Reference-only fitting uses receiver 0, while coordinates update one map at a time. Adapted spectral synchronization \cite{pachauri2013} selects pairwise relative maps, then uses a leading eigenspace and Hungarian rounding. Spectral-plus-coordinate refines the full RF objective. Exact fitting retains all pair-table scores; sequential fitting solves once, whereas alternation also proposes groups. All maps freeze before testing (Supplementary Section S15).

Five adaptation scenes give exact recovery of all 120 hidden entries for every scorer. Weighted, unweighted, and boosted F1 are 0.7786, 0.7997, and 0.8232, matching known-map references. These 120 entries belong to ten interventions, not 120 independent trials (Table~\ref{tab:deep_mapping}).

\begin{table}[!t]\begin{adjustwidth}{-\extralength}{0cm}
\caption{Missing-map estimators using five adaptation and 80 test scenes per replication. Three four-class maps are unknown; only the reference is known. Counts sum ten interventions; F1 intervals use ten replication means. Sequential exact, alternating exact, and known-map results coincide. All outputs satisfy receiver and gate constraints.}\label{tab:deep_mapping}
\small\setlength{\tabcolsep}{4pt}\begin{tabularx}{\linewidth}{@{}lYrr@{}}\toprule
Scorer & Map estimator & Correct / hidden & Event-pair F1 [95\% CI] \\\midrule
Weighted linear & Unresolved & 21/120 & 0.5316 [0.5119, 0.5514] \\
Weighted linear & Reference only & 113/120 & 0.7565 [0.7224, 0.7906] \\
Weighted linear & Spectral sync. & 114/120 & 0.7712 [0.7548, 0.7877] \\
Weighted linear & Spectral + coordinate & 114/120 & 0.7712 [0.7548, 0.7877] \\
Weighted linear & Coordinate & 70/120 & 0.7005 [0.6541, 0.7469] \\
Weighted linear & Exact block & 120/120 & 0.7786 [0.7718, 0.7853] \\
\addlinespace
Unweighted linear & Unresolved & 21/120 & 0.4695 [0.4452, 0.4937] \\
Unweighted linear & Reference only & 118/120 & 0.7912 [0.7713, 0.8110] \\
Unweighted linear & Spectral sync. & 120/120 & 0.7997 [0.7913, 0.8080] \\
Unweighted linear & Spectral + coordinate & 120/120 & 0.7997 [0.7913, 0.8080] \\
Unweighted linear & Coordinate & 66/120 & 0.7078 [0.6478, 0.7678] \\
Unweighted linear & Exact block & 120/120 & 0.7997 [0.7913, 0.8080] \\
\addlinespace
Boosted & Unresolved & 21/120 & 0.4458 [0.4230, 0.4686] \\
Boosted & Reference only & 114/120 & 0.7995 [0.7645, 0.8345] \\
Boosted & Spectral sync. & 114/120 & 0.8134 [0.7923, 0.8344] \\
Boosted & Spectral + coordinate & 114/120 & 0.8134 [0.7923, 0.8344] \\
Boosted & Coordinate & 69/120 & 0.7290 [0.6754, 0.7826] \\
Boosted & Exact block & 120/120 & 0.8232 [0.8183, 0.8282] \\
\addlinespace
\bottomrule\end{tabularx}

\end{adjustwidth}\end{table}

Exact-minus-spectral F1 is 0.0073 for weighted scoring with interval $[-0.0093,0.0239]$, and 0.0099 for boosting with $[-0.0125,0.0322]$. Unweighted results coincide, so no general superiority is established. At 40 scenes, reference-only, spectral, and exact fitting recover all maps, while coordinates can stall. Sequential and alternating maps coincide in all 90 scorer--budget--replication settings.

Two same-class conditions hide maps using five adaptation and 40 test scenes per replication. Exact fitting recovers every entry, yet at copy probability 0.65 the scorers give F1 0.4450, 0.4504, and 0.4537, equal to known-map references. Correct registration cannot distinguish nearby same-family events; the next test supplies all maps to isolate this ambiguity.

A five-scene weighted case has 160 grouped pairs and 13,824 states. Cached enumeration takes 2.44 ms versus 1413.08 ms for full reevaluation, with identical optima in three-run medians. A six-class, two-unknown kernel has 518,400 states, taking 377.79 ms for tables and 1.76 ms for search. These timings do not establish larger-vocabulary recovery; Supplementary Section S15 gives memory and allocation limits.

\subsection{Class evidence under harder descriptor conditions}\label{sec:deep_semantic}
Eight conditions use 400 new scenes each, with ten replications of 40 scenes. Four frozen scorers receive correct maps to isolate class-evidence quality. Semantic interventions preserve physical reports; dense conditions share a separate matched physical reference. Supplementary Section S15 specifies probability models, seeds, and intervals.

Nominal near-collisions have independent classes, with agreement probability $1/4$: 312 of 1303 copied-event parent pairs match in the new sample. Forcing copies to share the parent's class lowers weighted, unweighted, and boosted F1 from 0.7747, 0.7925, and 0.8226 to 0.7010, 0.7096, and 0.7279, respectively. Physical-only F1 remains 0.6773, attributing part of the semantic gain to class independence.

Increasing copy probability from 0.25 to 0.65 gives boosted F1 of 0.6335 with independent classes and 0.4482 with shared classes, versus physical-only F1 of 0.4460. The paired boosted-minus-physical gain is 0.0021 with interval $[-0.0113,0.0155]$; both linear intervals also include zero. Correct registration offers little event discrimination here.

\begin{table}[!t]\begin{adjustwidth}{-\extralength}{0cm}
\caption{Semantic stress tests: mean event-pair F1 across ten replication means, with 400 scenes per row and known maps. Semantic interventions preserve physical observations; dense rows share a separate physical reference. Error marginals match in law, though realized counts can differ. All outputs satisfy receiver and gate constraints. Supplementary Section S15 gives full intervals.}\label{tab:deep_semantic}
\small\setlength{\tabcolsep}{4pt}\begin{tabularx}{\linewidth}{@{}Yrrrr@{}}\toprule
Condition & Physical & Weighted linear & Unweighted linear & Boosted \\\midrule
Nominal independent classes & 0.6773 & 0.7747 & 0.7925 & 0.8226 \\
Same-class near collisions & 0.6773 & 0.7010 & 0.7096 & 0.7279 \\
Dense near collisions / independent classes & 0.4460 & 0.5831 & 0.6210 & 0.6335 \\
Dense near collisions / same class & 0.4460 & 0.4408 & 0.4471 & 0.4482 \\
Class prior 0.7/0.1/0.1/0.1 & 0.6773 & 0.7451 & 0.7606 & 0.7852 \\
Receiver-specific confusion and confidence & 0.6773 & 0.7601 & 0.7822 & 0.7992 \\
Independent errors / matched marginals & 0.6773 & 0.7714 & 0.7905 & 0.8184 \\
Event-shared errors / matched marginals & 0.6773 & 0.8334 & 0.8826 & 0.9081 \\
\bottomrule\end{tabularx}

\end{adjustwidth}\end{table}

Class imbalance and receiver-specific classifiers reduce gains (Table~\ref{tab:deep_semantic}). The latter retain error probability 0.12 but vary confusion destinations and confidence ranges, separating correct column registration from reliable recognition.

Sharing one error flag and wrong family across an event preserves marginal laws from an independent-error control. Weighted F1 rises from 0.7714 to 0.8334, a paired gain of 0.0620 with interval $[0.0532,0.0708]$. Agreement on a wrong family can preserve event compatibility, so event F1 does not certify class correctness or establish a universal effect of correlated errors.


\subsection{Risk under departures from the observation assumptions}\label{sec:risks_revision}
The first 20 tests per replication support nine interventions, two frozen scorers, and three decoders, totaling 10,800 outcomes. Clock equicorrelations 0.8 and $-0.25$ compare diagonal and matched declarations. Their small effects here do not establish robustness to arbitrary dependence.

Adding exponential delay with mean $8\,\mu$s to 30\% of genuine reports lowers gate recall from 0.9997 to 0.9353 and constrained F1 from 0.7822 to 0.7122. This is descriptor contamination, not measured NLOS validation.

The alarm test adds 12 false reports by copying three templates to all four receivers. Templates share time, carrier, bandwidth, duration, and canonical class vectors before independent Gaussian timing jitter of s.d. $0.2\,\mu$s and frequency jitter of 120 Hz. Native classes are re-encoded. The independent control uses separate templates with matched counts and marginal pool. Each false report has singleton truth; no common interfering waveform is generated.

Correlated false reports have clutter attachment 0.7821 versus 0.7061 in the independent control, yet F1 is also higher at 0.6023 versus 0.5105 (Table~\ref{tab:review_risks}). Both produce feasible false groups; dependence does not uniformly reduce F1.

With probability 0.2, a genuine report receives a same-receiver copy carrying its transmission label. The original stays unchanged; the copy gets independent Gaussian timing jitter of s.d. $0.15\,\mu$s and frequency jitter of 40 Hz, while bandwidth, duration, and class probabilities remain copied. These are detection records, not simulated waveform segments. Constrained F1 is 0.6009.

Let $n_{es}$ count reports of event $e$ at receiver $s$. The scene-level pair-recall ceiling under receiver exclusivity alone is
\[
 R_{\max}=\frac{\sum_e\sum_{s<t}\min(n_{es},n_{et})}
 {\sum_e\binom{\sum_s n_{es}}2}.
\]
The denominator includes same-receiver fragment pairs, which exclusive groups cannot recover. Each receiver pair permits at most $\min(n_{es},n_{et})$ links, attainable by layered grouping without gates. Empty-positive scenes receive ceiling zero. Figure~\ref{fig:review_risks}b averages scene ceilings within replication and then across ten replications. The fragment value 0.6529 bounds recall, not F1 (Supplementary Section S12).

\begin{table}[!t]\begin{adjustwidth}{-\extralength}{0cm}
\caption{Observation-assumption tests with ten replications of 20 scenes. Rx/gate entries are violation percentages; constrained zeros follow by construction. Clutter attachment counts false reports in non-singletons. The recall ceiling ignores gates and scores.}\label{tab:review_risks}
\small\setlength{\tabcolsep}{3pt}\begin{tabularx}{\linewidth}{@{}YYrrrr@{}}\toprule
Intervention & Event-pair F1 [95\% CI] & Gate recall & Clutter attachment & Recall ceiling & Rx / gate \\\midrule
Nominal subset & 0.7822 [0.7635, 0.8009] & 0.9997 & 0.0067 & 1.0000 & 0 / 0 \\
Positive corr.; diagonal declaration & 0.7816 [0.7618, 0.8014] & 0.9997 & 0.0067 & 1.0000 & 0 / 0 \\
Positive corr.; matched covariance & 0.7823 [0.7627, 0.8019] & 0.9997 & 0.0067 & 1.0000 & 0 / 0 \\
Negative corr.; diagonal declaration & 0.7796 [0.7612, 0.7979] & 0.9997 & 0.0057 & 1.0000 & 0 / 0 \\
Negative corr.; matched covariance & 0.7798 [0.7609, 0.7986] & 0.9997 & 0.0057 & 1.0000 & 0 / 0 \\
Extra positive delay & 0.7122 [0.6946, 0.7298] & 0.9353 & 0.0056 & 1.0000 & 0 / 0 \\
Independent false alarms; same count & 0.5105 [0.4993, 0.5217] & 0.9997 & 0.7061 & 1.0000 & 0 / 0 \\
Correlated false alarms & 0.6023 [0.5868, 0.6178] & 0.9997 & 0.7821 & 1.0000 & 0 / 0 \\
Fragmented reports & 0.6009 [0.5857, 0.6161] & 0.9998 & 0.0084 & 0.6529 & 0 / 0 \\
\bottomrule\end{tabularx}

\end{adjustwidth}\end{table}
\widefig{review_risks}{Assumption failures under frozen evidence. (a) Event F1 with 95\% replication intervals. (b) Gate recall and the exclusivity-only recall ceiling. Fragments can be unrepresentable despite admission, and correlated false reports can be feasible. Alarm controls have equal added-report counts.}{fig:review_risks}

\subsection{Medium-size optimization audit}\label{sec:medium_exact}
Report count alone selects the first two tests with at least 40 reports per replication. Nested random subsets of 20, 30, and 40 yield 60 MILPs under the frozen greedy objective. All 60 prove optimality within five-second limits at requested gap zero. Records retain incumbents, bounds, status, nodes, feasibility, and construction-plus-solve time.

Greedy mean objective losses are 0.18\%, 0.21\%, and 0.46\%, without implying a general bound. At 40 reports, restarts reduce loss to 0.11\%; greedy, restart, and exact F1 are 0.7974, 0.8098, and 0.8142 (Table~\ref{tab:review_exact}). Better objective search helps this sample but does not establish event truth.

\begin{table}[!t]\begin{adjustwidth}{-\extralength}{0cm}
\caption{Medium-size audit against proven optima. Times are medians in the review environment, including MILP construction. Every decoder is feasible; unproven incumbents are not labeled exact.}\label{tab:review_exact}
\small\setlength{\tabcolsep}{3pt}\begin{tabularx}{\linewidth}{@{}rrrrrrrY@{}}\toprule
$N$ & Proven optima & Greedy gap & Restart gap & Greedy F1 & Restart F1 & MILP F1 & Median ms: greedy / restart / MILP \\\midrule
20 & 20/20 & 0.2\% & 0.0\% & 0.7689 & 0.7750 & 0.7750 & 0.11 / 0.53 / 5.07 \\
30 & 20/20 & 0.2\% & 0.0\% & 0.7621 & 0.7649 & 0.7649 & 0.14 / 0.94 / 14.63 \\
40 & 20/20 & 0.5\% & 0.1\% & 0.7974 & 0.8098 & 0.8142 & 0.22 / 1.52 / 35.32 \\
\bottomrule\end{tabularx}

\end{adjustwidth}\end{table}

\paragraph{Fixed-registry processing time.}
Fresh report-path timing includes feature construction, scoring, and decoding for ten scenes at each of 40, 80, and 120 reports. Greedy medians are 0.64, 1.80, and 3.63 ms; triangle MILP medians are 34.56, 284.37, and 956.65 ms. All 30 solves prove optimality. The one-second cap excludes construction, which raises the maximum path to 1014.16 ms. These implementation timings assume no deployment deadline and exclude waveforms and offline map fitting. At 40 reports, greedy F1 is 0.8068 versus MILP F1 of 0.7781. Supplementary Section S15 records timings, solver status, and memory scope.

\subsection{Auxiliary public-recording metadata check}\label{sec:public_data}
\subsubsection{Source, inclusion rules, and independent capture blocks}
We use the publisher-selected six-file urban LoRaWAN sample \cite{loed2020,loeddataset}, released under CC BY 4.0 with recorded publisher MD5 and download SHA-256. Training uses 3 March 2019, 2 April 2019, and 2 May 2020; calibration uses 3 July 2020; testing uses 3 August and 2 September 2020. Recorded timestamps, including those preceding filename dates, are retained.

SHA-256 of the complete decoded payload defines the operational packet reference for labels and supervised fitting. Eligible records are CRC-valid unconfirmed uplinks with valid base64; message type comes from the payload header, avoiding the CSV's decimal-rendered binary labels. After exact deduplication, we exclude groups with repeated gateways, spans above 5 s, or identities shared across splits.

These rules remove obvious repetitions without independently proving one transmission. We therefore evaluate \emph{reference-packet association}, with split totals in Table~\ref{tab:loed_inclusion} and per-file exclusions in Supplementary Section S7.

\begin{table}[!t]\begin{adjustwidth}{-\extralength}{0cm}
\caption{LoED sample and inclusion counts. A visible reference packet has reports at two or more gateways. Eligible test reports come from two gateways. The original dataset covers nine gateways. This experiment does not use all nine.}\label{tab:loed_inclusion}
\small\renewcommand{\arraystretch}{1.12}\begin{tabular*}{\linewidth}{@{\extracolsep{\fill}}lrrrr@{}}\toprule
Split & Files & Raw reports & Eligible reports & Visible reference packets \\\midrule
Training & 3 & 260,231 & 39,447 & 257 \\
Calibration & 1 & 25,659 & 17,252 & 434 \\
Test & 2 & 40,980 & 23,677 & 1,417 \\
\bottomrule\end{tabular*}

\end{adjustwidth}\end{table}

Repeated payloads at individual gateways explain many 2019 exclusions. Results therefore concern identifiable receiver-exclusive references, not all traffic; CRC-invalid reports have no independent clutter labels. Reference hashes are stored separately, and scoring excludes addresses, frame counters, payloads and hashes, ports, and row indices.

\subsubsection{Measured-input adaptation and frozen parameters}
LoED has no waveform-derived uncertainty or class posteriors, so the reduced model uses reception time $t_i$, channel $F_i$, and decoded spreading factor $L_i$. Clock corrections $\widehat b_s$ use matching training references.

Gateway pairs need ten references for clock fitting. Median differences define a weighted least-squares system whose minimum-norm solution fixes the common time origin; unsupported corrections remain zero. The supported pair median is 0.999985 s, without test-identity fitting.

Corrected time is $t_i'=t_i-\widehat b_{s_i}$, with tolerance $G=\max(1\,\mathrm{ms},Q_{0.99})=3.02517\,\mathrm{ms}$. Here $Q_{0.99}$ is the 99th percentile of absolute calibration-reference residuals. This empirical tolerance does not carry the Gaussian guarantee of Proposition~\ref{prop:gate}. The features are
\begin{equation}
 \psi_{ij}=\left[\log\!\left(1+\frac{|t_i'-t_j'|}{G}\right),\;
 \ind[F_i\ne F_j],\;\ind[L_i\ne L_j]\right].\label{eq:loed_features}
\end{equation}
Timing scoring uses the first coordinate; metadata scoring uses all three. Inputs are standardized and logistic coefficients fitted on training data, with temperature fitted on separate calibration data. Admission requires $|t_i'-t_j'|\leq G$.

Reported channel frequency is not oscillator offset, and decoded spreading factor is not a probabilistic recognition output. We synthesize no geometry, RF uncertainty, or class posterior.

Corrected timestamps define fixed 60 s windows without identity-based reassignment. Each scorer uses connected components, pairwise assignment followed by union, and constrained decoding.

Evaluation retains singletons and counts window-split reference pairs as false negatives over each complete capture. We report captures separately and pool TP, FP, and FN. Two captures do not support deployment-population intervals based on individual pairs.

\subsubsection{Measured outcomes and interpretation}
The 23,677 reports contain 22,260 reference packets, including 1417 observed at both gateways. Metadata constraints produce 1411 true-positive, one false-positive, and six false-negative pairs: precision 0.9993, recall 0.9958, and F1 0.9975. On 3 August, TP is 900, FP 0, and FN 3 (F1 0.9983); on 2 September, TP is 511, FP 1, and FN 3 (F1 0.9961). Assignment and components give the same partition.

Timing-only components produce seven false positives and one receiver conflict. Assignment and constraints each produce five false positives with F1 0.9961 (Table~\ref{tab:loed_results}).

\begin{table}[!t]\begin{adjustwidth}{-\extralength}{0cm}
\caption{Public-recording results. Values are reference-packet F1. Rx counts clusters with repeated receivers. Within each scorer, all methods use the same observations, scores, gates, and capture windows. We pool pair counts before calculating F1. We do not average rounded scores.}\label{tab:loed_results}
\small\renewcommand{\arraystretch}{1.12}\begin{tabular*}{\linewidth}{@{\extracolsep{\fill}}llrrrr@{}}\toprule
Evidence & Decoder & 3 Aug. F1 & 2 Sep. F1 & Pooled F1 & Rx \\\midrule
Timing & Components & 0.9972 & 0.9922 & 0.9954 & 1 \\
Timing & Assignment & 0.9983 & 0.9922 & 0.9961 & 0 \\
Timing & Constrained & 0.9983 & 0.9922 & 0.9961 & 0 \\
Metadata & Components & 0.9983 & 0.9961 & 0.9975 & 0 \\
Metadata & Assignment & 0.9983 & 0.9961 & 0.9975 & 0 \\
Metadata & Constrained & 0.9983 & 0.9961 & 0.9975 & 0 \\
\bottomrule\end{tabular*}

\end{adjustwidth}\end{table}

Shared decoded meanings support a reproducible metadata-to-packet procedure, but metadata already resolves nearly all ambiguity. The two-gateway check has no measured class posteriors or source-specific protocol vocabulary and offers little room to distinguish decoder accuracy.

Assignment is already receiver-exclusive with two gateways, leaving higher-order grouping untested. Metadata F1 0.9975 does not validate multi-protocol robustness or RF uncertainty. Section~\ref{sec:public_sensitivity} tests the supervised preprocessing behind this result.

\subsection{Sensitivity of the auxiliary public check}\label{sec:public_sensitivity}
The public sensitivity audit fixes all 23,677 test IDs while varying training, calibration, clock correction, or quantile (Table~\ref{tab:review_public}). Original F1 is 0.9975, unchanged with May 2020 training alone. The two 2019 blocks provide one positive pair and no clock pair meeting ten-reference support, giving zero F1. Two settings lack usable training classes or calibration matches and are reported NE.

Without clock correction, metadata F1 is 0.9457. Quantile 0.95 gives F1 0.9933; quantile 0.999 expands the gate to 3.31 s and gives 0.9905. More candidates need not improve association. The original 0.99 rule stays fixed without test-F1 selection.

\begin{table}[!t]\begin{adjustwidth}{-\extralength}{0cm}
\caption{Public preprocessing sensitivity with unchanged test reports. Train positives count reference pairs available to scorer fitting. Clock corrections require at least ten matches for a receiver pair. $G$ is the empirical time gate. F1 is reference-packet F1. NE means the required fit cannot be estimated.}\label{tab:review_public}
\small\setlength{\tabcolsep}{3pt}\begin{tabularx}{\linewidth}{@{}Yrrrrr@{}}\toprule
Preprocessing setting & Train positives & $G$ (ms) & Aug. F1 & Sep. F1 & Pooled F1 \\\midrule
Original 3/1 split & 257 & 3.025 & 0.9983 & 0.9961 & 0.9975 \\
2019 training blocks only & 1 & 1003.010 & 0.0000 & 0.0000 & 0.0000 \\
May training only & 256 & 3.025 & 0.9983 & 0.9961 & 0.9975 \\
April training only & 0 & -- & NE & NE & NE \\
March train / April calibrate & -- & -- & NE & NE & NE \\
2019 train / May calibrate & 1 & 2000.887 & 0.0000 & 0.0000 & 0.0000 \\
Clock correction disabled & 257 & 1003.010 & 0.9603 & 0.9202 & 0.9457 \\
Calibration quantile 0.95 & 257 & 1.000 & 0.9956 & 0.9892 & 0.9933 \\
Calibration quantile 0.999 & 257 & 3307.621 & 0.9923 & 0.9875 & 0.9905 \\
\bottomrule\end{tabularx}

\end{adjustwidth}\end{table}

\subsection{Public-reference tails and clock supervision}\label{sec:deep_loed}
Of July's 434 reference pairs, four have corrected differences above two seconds despite matching payloads and counters; their reported channels differ. Five May pairs and two September pairs show the same pattern, while all 2096 other 2020 pairs have matching channels. Original CSV checks confirm these fields. A single timestamp without an independent transmission identifier cannot distinguish repetitions from metadata errors; all pairs remain evaluated.

July's 0.99 quantile interpolates between 2.120 and 3.471 ms; the 0.999 quantile lies between 2.136977 and 4.201604 s, explaining the jump to 3.308 s. Channel disagreement further limits the payload reference's physical meaning, precluding interpretation as clock drift alone.

A reference-free training lag histogram uses time, channel, spreading factor, and gateway ID with fixed support and peak-separation rules. Refinement contains 251 candidates with median 0.999985 s, matching the supervised offset numerically. Both fits give identical partitions on 23,677 reports, with TP 1411, FP 1, FN 6, and F1 0.9975. Population selection, scoring, and gate calibration remain reference-assisted (Supplementary Section S15).

The 2019-only fit has one positive pair from gateways absent in the 2020 tests, and April has no cross-gateway reference pair. Zero F1 combines sparse supervision and unsupported calibration rather than isolating drift. These retrospective checks diagnose preprocessing and reference construction.

\FloatBarrier

